\documentclass[../../main.tex]{subfiles}

\begin{document}

\section{非负可测函数的积分}\label{section:ch18-nonnegative-integral}

令 $(X,\mathcal{M},\mu)$ 为测度空间. 我们首先定义 $X$ 上的非负简单函数的积分. 令 $\psi$ 为非负简单函数, 它表示为
\begin{align*}
\psi=\sum_{k=1}^{n}c_{k}\cdot\chi_{E_{k}},
\end{align*}
其中 $c_{k}\geqslant 0$ 是实数且 $E_{k}$ 是可测集. $\psi$ 在 $X$ 上关于 $\mu$ 的积分, 记为 $\int_{X}\psi\,d\mu$, 定义为
\begin{align}
\int_{X}\psi\,d\mu=\sum_{k=1}^{n}c_{k}\cdot\mu(E_{k}).\label{align:ch18-simple-integral}
\end{align}
由于 $\psi$ 有一个规范表示(即不交的可测集族上取值不同的常数), 且在规范表示中系数是 $\psi$ 的取值, 积分的定义是恰当的. 注意若某个 $c_{k}=0$, 则 $\mu(E_{k})$ 必须为使得 $0\cdot\infty=0$ 才有意义.

对 $X$ 上的非负可测函数 $f$, $f$ 在 $X$ 上关于 $\mu$ 的积分, 记为 $\int_{X}f\,d\mu$, 定义为
\begin{align}
\int_{X}f\,d\mu=\sup\left\{\int_{X}\psi\,d\mu\,\middle|\,0\leqslant\psi\leqslant f,\,\psi\text{ 是简单函数}\right\}.\label{align:ch18-nonneg-integral-def}
\end{align}
对可测集 $E$ 和 $X$ 上的非负可测函数 $f$, $f$ 在 $E$ 上的积分定义为
\begin{align*}
\int_{E}f\,d\mu=\int_{X}f\cdot\chi_{E}\,d\mu.
\end{align*}

以下命题的证明从积分定义与简单函数的性质直接得到.

\begin{proposition}\label{proposition:ch18-integral-basic}
令 $(X,\mathcal{M},\mu)$ 为测度空间.
\begin{description}
\item[（线性）] 对非负简单函数 $\psi$ 和 $\varphi$ 以及非负实数 $\alpha$ 和 $\beta$,
\begin{align*}
\int_{X}(\alpha\psi+\beta\varphi)\,d\mu=\alpha\int_{X}\psi\,d\mu+\beta\int_{X}\varphi\,d\mu.
\end{align*}
\item[（单调性）] 若 $0\leqslant f\leqslant g$ 在 $X$ 上 a.e. 成立, 则
\begin{align*}
\int_{X}f\,d\mu\leqslant\int_{X}g\,d\mu.
\end{align*}
\item[（区域上的可加性）] 若 $E_{1}$ 和 $E_{2}$ 是 $X$ 的不交可测子集, 则对非负可测函数 $f$,
\begin{align*}
\int_{E_{1}\cup E_{2}}f\,d\mu=\int_{E_{1}}f\,d\mu+\int_{E_{2}}f\,d\mu.
\end{align*}
\item[（零测集上的积分）] 若 $E$ 是可测集且 $\mu(E)=0$, 则对非负可测函数 $f$,
\begin{align*}
\int_{E}f\,d\mu=0.
\end{align*}
\end{description}
\end{proposition}

我们现在证明非负可测函数的积分理论的中心结论——Fatou 引理.

\begin{lemma}[Fatou]\label{lemma:ch18-fatou}
令 $(X,\mathcal{M},\mu)$ 为测度空间, 而 $\{f_{n}\}$ 是 $X$ 上的非负可测函数序列. 则
\begin{align}
\int_{X}\liminf_{n\to\infty}f_{n}\,d\mu\leqslant\liminf_{n\to\infty}\int_{X}f_{n}\,d\mu.\label{align:ch18-fatou}
\end{align}
\end{lemma}

\begin{proof}
定义 $f=\liminf_{n\to\infty}f_{n}$. 由于每个 $f_{n}$ 是非负可测的, $f$ 也是非负可测的. 令 $\psi$ 为满足 $0\leqslant\psi\leqslant f$ 的简单函数. 我们将证明
\begin{align}
\int_{X}\psi\,d\mu\leqslant\liminf_{n\to\infty}\int_{X}f_{n}\,d\mu.\label{align:ch18-fatou-proof1}
\end{align}
由于 $\int_{X}f\,d\mu$ 是所有这样的 $\psi$ 的积分的上确界, 这就推出\eqref{align:ch18-fatou}.

将 $\psi$ 表示为 $\psi=\sum_{k=1}^{m}c_{k}\cdot\chi_{E_{k}}$, 其中 $\{E_{k}\}_{k=1}^{m}$ 是 $X$ 的不交可测集族且每个 $c_{k}>0$. 若 $\int_{X}\psi\,d\mu=\infty$, 则存在某个 $k$ 使得 $\mu(E_{k})=\infty$ 且 $c_{k}>0$. 由于 $\psi\leqslant f$, 在 $E_{k}$ 上 $f\geqslant c_{k}$. 因此, 对每个自然数 $n$,
\begin{align*}
f=\liminf_{j\to\infty}f_{j}\geqslant c_{k}\text{ 在 }E_{k}\text{ 上}.
\end{align*}
令 $\alpha$ 为满足 $0<\alpha<1$ 的实数. 对每个自然数 $n$, 定义
\begin{align*}
E_{n}=\{x\in E_{k}\mid f_{j}(x)\geqslant\alpha c_{k}\text{ 对所有 }j\geqslant n\}.
\end{align*}
由于 $f_{j}$ 是可测的, 每个 $E_{n}$ 是可测的. 此外, $\{E_{n}\}$ 是上升的序列且 $E_{k}=\bigcup_{n=1}^{\infty}E_{n}$, 因为在 $E_{k}$ 上 $\liminf_{j\to\infty}f_{j}\geqslant c_{k}$. 因此, 对每个 $n$,
\begin{align*}
\int_{X}f_{j}\,d\mu\geqslant\int_{E_{n}}f_{j}\,d\mu\geqslant\alpha c_{k}\cdot\mu(E_{n})\text{ 对所有 }j\geqslant n.
\end{align*}
因此, $\liminf_{j\to\infty}\int_{X}f_{j}\,d\mu\geqslant\alpha c_{k}\cdot\mu(E_{n})$. 由于 $\mu(E_{k})=\infty$ 且 $\{E_{n}\}$ 上升到 $E_{k}$, 根据测度的连续性, $\lim_{n\to\infty}\mu(E_{n})=\mu(E_{k})=\infty$. 因此 $\liminf_{j\to\infty}\int_{X}f_{j}\,d\mu=\infty$, 而\eqref{align:ch18-fatou-proof1}成立.

现在假设 $\int_{X}\psi\,d\mu<\infty$. 令 $\alpha$ 为满足 $0<\alpha<1$ 的实数. 对每个自然数 $n$ 和每个 $k$, 定义
\begin{align*}
E_{n,k}=\{x\in E_{k}\mid f_{j}(x)\geqslant\alpha c_{k}\text{ 对所有 }j\geqslant n\}.
\end{align*}
则 $\{E_{n,k}\}_{n=1}^{\infty}$ 是上升的可测集序列且 $E_{k}=\bigcup_{n=1}^{\infty}E_{n,k}$. 对每个 $n$, 定义
\begin{align*}
\psi_{n}=\sum_{k=1}^{m}\alpha c_{k}\cdot\chi_{E_{n,k}}.
\end{align*}
则 $\psi_{n}$ 是简单函数且对所有 $j\geqslant n$, $\psi_{n}\leqslant f_{j}$ 在 $X$ 上成立. 因此, 对所有 $j\geqslant n$,
\begin{align*}
\int_{X}f_{j}\,d\mu\geqslant\int_{X}\psi_{n}\,d\mu=\sum_{k=1}^{m}\alpha c_{k}\cdot\mu(E_{n,k}).
\end{align*}
取 $\liminf_{j\to\infty}$ 得到
\begin{align*}
\liminf_{j\to\infty}\int_{X}f_{j}\,d\mu\geqslant\sum_{k=1}^{m}\alpha c_{k}\cdot\mu(E_{n,k}).
\end{align*}
现在令 $n\to\infty$. 根据测度的连续性, $\lim_{n\to\infty}\mu(E_{n,k})=\mu(E_{k})$. 因此
\begin{align*}
\liminf_{j\to\infty}\int_{X}f_{j}\,d\mu\geqslant\alpha\sum_{k=1}^{m}c_{k}\cdot\mu(E_{k})=\alpha\int_{X}\psi\,d\mu.
\end{align*}
最后, 令 $\alpha\to 1^{-}$ 得到\eqref{align:ch18-fatou-proof1}.
\end{proof}

从 Fatou 引理我们立即得到两个关键的收敛定理.

\begin{theorem}[单调收敛定理]\label{theorem:ch18-monotone-convergence}
令 $(X,\mathcal{M},\mu)$ 为测度空间, 而 $\{f_{n}\}$ 是 $X$ 上的非负可测函数序列, 满足在 $X$ 上 $0\leqslant f_{1}\leqslant f_{2}\leqslant\cdots$ 且 $\lim_{n\to\infty}f_{n}=f$ 在 $X$ 上 a.e. 成立. 则
\begin{align}
\int_{X}f\,d\mu=\lim_{n\to\infty}\int_{X}f_{n}\,d\mu.\label{align:ch18-mct}
\end{align}
\end{theorem}

\begin{proof}
由于 $\{f_{n}\}$ 是上升的, $f=\lim_{n\to\infty}f_{n}=\liminf_{n\to\infty}f_{n}$. 根据 Fatou 引理,
\begin{align*}
\int_{X}f\,d\mu=\int_{X}\liminf_{n\to\infty}f_{n}\,d\mu\leqslant\liminf_{n\to\infty}\int_{X}f_{n}\,d\mu.
\end{align*}
另一方面, 由于对所有 $n$, $f_{n}\leqslant f$ a.e., 根据积分的单调性, $\int_{X}f_{n}\,d\mu\leqslant\int_{X}f\,d\mu$. 因此 $\limsup_{n\to\infty}\int_{X}f_{n}\,d\mu\leqslant\int_{X}f\,d\mu$. 将两个不等式联立得到\eqref{align:ch18-mct}.
\end{proof}

\begin{corollary}[Beppo Levi]\label{corollary:ch18-beppo-levi}
令 $(X,\mathcal{M},\mu)$ 为测度空间, 而 $\{f_{n}\}$ 是 $X$ 上的非负可测函数序列, 使得在 $X$ 上 $\sum_{n=1}^{\infty}f_{n}$ 收敛. 则
\begin{align}
\int_{X}\sum_{n=1}^{\infty}f_{n}\,d\mu=\sum_{n=1}^{\infty}\int_{X}f_{n}\,d\mu.\label{align:ch18-beppo-levi}
\end{align}
\end{corollary}

\begin{proof}
定义 $g_{n}=\sum_{k=1}^{n}f_{k}$. 则 $\{g_{n}\}$ 是上升的非负可测函数序列且 $\lim_{n\to\infty}g_{n}=\sum_{k=1}^{\infty}f_{k}$. 根据单调收敛定理,
\begin{align*}
\int_{X}\sum_{k=1}^{\infty}f_{k}\,d\mu=\lim_{n\to\infty}\int_{X}g_{n}\,d\mu.
\end{align*}
另一方面, 根据简单函数积分的线性性(对非负简单函数归纳), $\int_{X}g_{n}\,d\mu=\sum_{k=1}^{n}\int_{X}f_{k}\,d\mu$. 因此
\begin{align*}
\int_{X}\sum_{k=1}^{\infty}f_{k}\,d\mu=\lim_{n\to\infty}\sum_{k=1}^{n}\int_{X}f_{k}\,d\mu=\sum_{k=1}^{\infty}\int_{X}f_{k}\,d\mu.
\end{align*}
\end{proof}

Beppo Levi 引理的一个直接推论是积分在区域上的可数可加性.

\begin{corollary}\label{corollary:ch18-countable-additivity-integral}
令 $(X,\mathcal{M},\mu)$ 为测度空间, $E$ 是 $X$ 的可测子集, 而 $\{E_{k}\}_{k=1}^{\infty}$ 是 $E$ 的不交可测子集的可数族, 满足 $E=\bigcup_{k=1}^{\infty}E_{k}$. 则对 $X$ 上的非负可测函数 $f$,
\begin{align*}
\int_{E}f\,d\mu=\sum_{k=1}^{\infty}\int_{E_{k}}f\,d\mu.
\end{align*}
\end{corollary}

\begin{proof}
对每个 $k$, 令 $f_{k}=f\cdot\chi_{E_{k}}$. 则每个 $f_{k}$ 是非负可测的且 $f\cdot\chi_{E}=\sum_{k=1}^{\infty}f_{k}$. 根据 Beppo Levi 引理,
\begin{align*}
\int_{E}f\,d\mu=\int_{X}f\cdot\chi_{E}\,d\mu=\sum_{k=1}^{\infty}\int_{X}f_{k}\,d\mu=\sum_{k=1}^{\infty}\int_{E_{k}}f\,d\mu.
\end{align*}
\end{proof}

我们已经对非负简单函数证明了积分的线性性. 以下定理将线性性延拓到所有非负可测函数.

\begin{theorem}\label{theorem:ch18-nonneg-linearity}
令 $(X,\mathcal{M},\mu)$ 为测度空间, 而 $f$ 和 $g$ 是 $X$ 上的非负可测函数. 则对非负实数 $\alpha$ 和 $\beta$,
\begin{align*}
\int_{X}(\alpha f+\beta g)\,d\mu=\alpha\int_{X}f\,d\mu+\beta\int_{X}g\,d\mu.
\end{align*}
\end{theorem}

\begin{proof}
根据简单逼近定理(定理\cref{theorem:ch18-simple-approximation-thm}), 存在上升的非负简单函数序列 $\{\psi_{n}\}$ 和 $\{\varphi_{n}\}$, 使得在 $X$ 上逐点地 $\psi_{n}\to f$ 且 $\varphi_{n}\to g$. 则 $\{\alpha\psi_{n}+\beta\varphi_{n}\}$ 是上升的非负简单函数序列, 逐点收敛于 $\alpha f+\beta g$. 根据单调收敛定理,
\begin{align*}
\int_{X}(\alpha f+\beta g)\,d\mu=\lim_{n\to\infty}\int_{X}(\alpha\psi_{n}+\beta\varphi_{n})\,d\mu.
\end{align*}
但对简单函数, 线性性已被证明. 因此
\begin{align*}
\int_{X}(\alpha\psi_{n}+\beta\varphi_{n})\,d\mu=\alpha\int_{X}\psi_{n}\,d\mu+\beta\int_{X}\varphi_{n}\,d\mu.
\end{align*}
再由单调收敛定理, $\lim_{n\to\infty}\int_{X}\psi_{n}\,d\mu=\int_{X}f\,d\mu$ 且 $\lim_{n\to\infty}\int_{X}\varphi_{n}\,d\mu=\int_{X}g\,d\mu$. 因此定理得证.
\end{proof}

\begin{exercise}\label{exercise:ch18-17}
令 $(X,\mathcal{M},\mu)$ 为测度空间, 而 $f$ 是 $X$ 上的非负可测函数, 满足 $\int_{X}f\,d\mu=0$. 证明 $f=0$ a.e.
\end{exercise}

\begin{exercise}\label{exercise:ch18-18}
令 $(X,\mathcal{M},\mu)$ 为测度空间. 定义集函数 $\nu(E)=\int_{E}f\,d\mu$, 其中 $f$ 是 $X$ 上的非负可测函数. 证明 $\nu$ 是 $\mathcal{M}$ 上的测度.
\end{exercise}

\begin{exercise}\label{exercise:ch18-19}
令 $\{f_{n}\}$ 为 $X$ 上的非负可测函数序列. 证明
\begin{align*}
\int_{X}\sum_{n=1}^{\infty}f_{n}\,d\mu=\sum_{n=1}^{\infty}\int_{X}f_{n}\,d\mu.
\end{align*}
（这是 Beppo Levi 引理的另一种表述.）
\end{exercise}

\begin{exercise}\label{exercise:ch18-20}
令 $(X,\mathcal{M},\mu)$ 为测度空间, 而 $f$ 是 $X$ 上的非负可测函数. 定义 $E_{n}=\{x\in X\mid f(x)\geqslant n\}$. 证明
\begin{align*}
\int_{X}f\,d\mu\geqslant\sum_{n=1}^{\infty}\mu(E_{n}).
\end{align*}
若 $f$ 是有界的, 证明等号成立.
\end{exercise}

\begin{exercise}\label{exercise:ch18-21}
证明 Levi 单收敛定理的逆否命题: 若 $\{f_{n}\}$ 是 $X$ 上的非负可测函数序列且 $\lim_{n\to\infty}\int_{X}f_{n}\,d\mu<\infty$, 则 $\{f_{n}\}$ 在 $X$ 上 a.e. 有限.
\end{exercise}

\begin{exercise}\label{exercise:ch18-22}
令 $(X,\mathcal{M},\mu)$ 为测度空间, 而 $\{f_{n}\}$ 是 $X$ 上的非负可测函数序列, 满足在 $X$ 上 a.e. $f_{n}\to f$. 证明
\begin{align*}
\int_{X}f\,d\mu\leqslant\liminf_{n\to\infty}\int_{X}f_{n}\,d\mu.
\end{align*}
（这是 Fatou 引理.）
\end{exercise}

\end{document}